\textbf{Phần 1. Trắc nghiệm lựa chọn}
\setcounterpageref{ex}{0}

\begin{ex} %[2D2N12]
Trong bảng trừ (qua 10), kết quả của phép tính $11 - 9$ là
\choice
{\True $2$}
{$1$}
{$3$}
{$4$}
\loigiai{
$11 - 9 = 2$.\\
\textbf{Đáp án của câu hỏi là A}
}
\end{ex}

\begin{ex} %[2D2N12]
Phép tính nào sau đây có kết quả bằng 8?
\choice
{$14 - 7$}
{$13 - 6$}
{$15 - 8$}
{\True $14 - 6$}
\loigiai{
$14 - 6 = 8$.\\
\textbf{Đáp án của câu hỏi là D}
}
\end{ex}

\begin{ex} %[2D2H12]
Phép tính nào sau đây có kết quả bằng 7?
\choice
{\True $15 - 8$}
{$15 - 7$}
{$14 - 6$}
{$13 - 5$}
\loigiai{
$15 - 8 = 7$.\\
\textbf{Đáp án của câu hỏi là A}
}
\end{ex}

\begin{ex} %[2D2H12]
Phép tính nào sau đây có kết quả lớn nhất?
\choice
{$12 - 5$}
{$13 - 6$}
{\True $17 - 8$}
{$14 - 7$}
\loigiai{
$17 - 8 = 9$ (lớn nhất).\\
\textbf{Đáp án của câu hỏi là C}
}
\end{ex}

\begin{ex} %[2D2H12]
Phép tính nào sau đây có kết quả bé nhất?
\choice
{$13 - 8$}
{$12 - 7$}
{$14 - 8$}
{\True $11 - 7$}
\loigiai{
$11 - 7 = 4$ (bé nhất).\\
\textbf{Đáp án của câu hỏi là D}
}
\end{ex}

\begin{ex} %[2D2H12]
Điền số thích hợp: $13 - \dots = 9$.
\choice
{$3$}
{\True $4$}
{$5$}
{$6$}
\loigiai{
$13 - 4 = 9$.\\
\textbf{Đáp án của câu hỏi là B}
}
\end{ex}

\begin{ex} %[2D2V12]
Có bao nhiêu phép trừ dạng $11 - a$ trong bảng trừ? ($a$ từ 2 đến 9)
\choice
{$7$ phép tính}
{$9$ phép tính}
{\True $8$ phép tính}
{$10$ phép tính}
\loigiai{
Từ $11 - 2$ đến $11 - 9$ có $9 - 2 + 1 = 8$ phép tính.\\
\textbf{Đáp án của câu hỏi là C}
}
\end{ex}

\begin{ex} %[2D2V12]
Tính: $16 - 7 - 3 = \dots$
\choice
{\True $6$}
{$5$}
{$7$}
{$8$}
\loigiai{
$9 - 3 = 6$.\\
\textbf{Đáp án của câu hỏi là A}
}
\end{ex}

\begin{ex} %[2D2V12]
Điền dấu thích hợp: $15 - 7 \dots 16 - 8$.
\choice
{$>$}
{$<$}
{Không so sánh được}
{\True $=$}
\loigiai{
$15 - 7 = 8$ và $16 - 8 = 8$. Vậy $15 - 7 = 16 - 8$.\\
\textbf{Đáp án của câu hỏi là D}
}
\end{ex}

\begin{ex} %[2D2V12]
Một hộp có 18 viên bi, lấy ra 9 viên bi. Trong hộp còn lại số viên bi là
\choice
{$8$ viên}
{\True $9$ viên}
{$10$ viên}
{$11$ viên}
\loigiai{
$18 - 9 = 9$ viên.\\
\textbf{Đáp án của câu hỏi là B}
}
\end{ex}

\begin{ex} %[2D2V12]
Tìm $y$ biết: $14 - y = 8$.
\choice
{$5$}
{$7$}
{\True $6$}
{$8$}
\loigiai{
$y = 14 - 8 = 6$.\\
\textbf{Đáp án của câu hỏi là C}
}
\end{ex}

\begin{ex} %[2D2C12]
Hiệu của số lớn nhất có hai chữ số và số bé nhất có hai chữ số là bao nhiêu?
\choice
{$88$}
{\True $89$}
{$90$}
{$80$}
\loigiai{
$99 - 10 = 89$.\\
\textbf{Đáp án của câu hỏi là B}
}
\end{ex}

\textbf{Phần 2. Trắc nghiệm đúng sai}
\setcounterpageref{ex}{0}

\begin{ex} %[2D2H12]
Xét tính đúng sai của các phép trừ trong bảng trừ:
\choiceTF
{\True $13 - 6 = 7$.}
{\True $15 - 9 = 6$.}
{$17 - 9 = 9$.}
{\True $17 - 9 = 8$.}
\loigiai{
\begin{itemize}
    \item $13 - 6 = 7$: Đúng.
    \item $15 - 9 = 6$: Đúng.
    \item $17 - 9 = 8 \neq 9$. Ý 3 Sai, Ý 4 Đúng.
\end{itemize}
\begin{itemize}
    \item a) \textbf{Đúng}.
    \item b) \textbf{Đúng}.
    \item c) \textbf{Sai}.
    \item d) \textbf{Đúng}.
\end{itemize}
}
\end{ex}

\begin{ex} %[2D2V12]
Cho 3 phép tính: $11 - 6, 12 - 7, 13 - 8$. Xét tính đúng sai:
\choiceTF
{\True Cả 3 phép tính đều có kết quả bằng 5.}
{\True $11 - 6 = 5$.}
{Phép tính $13 - 8$ có kết quả lớn nhất.}
{\True Ba phép tính này có kết quả bằng nhau.}
\loigiai{
\begin{itemize}
    \item Cả 3 đều bằng 5: Đúng.
    \item $11 - 6 = 5$: Đúng.
    \item Chúng bằng nhau (không có lớn nhất). Ý 3 Sai, Ý 4 Đúng.
\end{itemize}
\begin{itemize}
    \item a) \textbf{Đúng}.
    \item b) \textbf{Đúng}.
    \item c) \textbf{Sai}.
    \item d) \textbf{Đúng}.
\end{itemize}
}
\end{ex}

\textbf{Phần 3. Trả lời ngắn}
\setcounterpageref{ex}{0}

\begin{ex} %[2D2H12]
Tính: $12 - 9$.
\shortans{3}
\loigiai{
\textbf{Đáp số:} $3$\\
$12 - 9 = 3$.
}
\end{ex}

\begin{ex} %[2D2H12]
Tính: $14 - 7$.
\shortans{7}
\loigiai{
\textbf{Đáp số:} $7$\\
$14 - 7 = 7$.
}
\end{ex}

\begin{ex} %[2D2V12]
Tính giá trị biểu thức: $15 - 8 + 4$.
\shortans{11}
\loigiai{
\textbf{Đáp số:} $11$\\
$7 + 4 = 11$.
}
\end{ex}

\begin{ex} %[2D2V12]
Điền số thích hợp: $18 - ? = 9$.
\shortans{9}
\loigiai{
\textbf{Đáp số:} $9$\\
$18 - 9 = 9$.
}
\end{ex}

\begin{ex} %[2D2V12]
Một cây có 16 quả bưởi, đã hái đi 8 quả. Hỏi trên cây còn lại bao nhiêu quả bưởi?
\shortans{8}
\loigiai{
\textbf{Đáp số:} $8$\\
$16 - 8 = 8$ quả.
}
\end{ex}

\begin{ex} %[2D2C12]
Tìm số thích hợp điền vào dấu ?: $13 - ? = 5$.
\shortans{8}
\loigiai{
\textbf{Đáp số:} $8$\\
$13 - 5 = 8$.
}
\end{ex}
